\documentclass[11pt,a4paper]{article}

\usepackage[utf8]{inputenc}

\usepackage[T1]{fontenc}

\usepackage[polish]{babel}

\usepackage{lmodern}

\usepackage{geometry}

\geometry{margin=27mm}

\usepackage[hidelinks]{hyperref}

\title{Nasilenie objawów pourazowych a deklarowane korzystanie z generatywnej AI do wsparcia emocjonalnego: badanie przekrojowe}

\author{Mateusz Pomianek}

\date{Szkic naukowy, wersja 1.0.0 | 28 listopada 2026}

\begin{document}

\maketitle

\begin{abstract}

Tematem pracy jest Nasilenie objawów pourazowych a deklarowane korzystanie z generatywnej AI do wsparcia emocjonalnego: badanie przekrojowe. Tekst systematyzuje uzasadnienie, hipotezy i sposób weryfikacji zagadnienia. Proponowany wskaźnik: Asocjacyjny OR z dwustronnym 95\% CI na 1 SD PCL-5 wobec binarnego Y\_E w zdefiniowanej populacji A\_X, zgodnie z SAP pracy. Nie przedstawia się niezweryfikowanych wyników ani fikcyjnych cytowań.

\end{abstract}

\section{Wprowadzenie}

PCL-5 mierzy nasilenie objawów pourazowych (Blevins et al., 2015). Badanie ustrukturyzowanego agenta tekstowego pokazuje precedens użycia narzędzia konwersacyjnego w obszarze zdrowia psychicznego (Fitzpatrick et al., 2017). Dotyczy odrębnej technologii i populacji, więc stanowi wyłącznie kontekst dla hipotezy o generatywnej AI.

\section{Cel i pytania badawcze}

Hipotezy robocze wymagają zatwierdzenia przed rozpoczęciem właściwej weryfikacji.

\paragraph{H1.} W zdefiniowanej populacji analitycznej $A_X$ wzrost wyniku PCL-5 o jedno odchylenie standardowe będzie związany z większymi szansami deklarowanego użycia generatywnej AI do wsparcia emocjonalnego, czyli OR dla $Y_E$ będzie większy od 1.

\section{Materiał i metody}

Pierwszorzędowy wskaźnik: iloraz szans deklarowanego użycia na jedno odchylenie standardowe wyniku PCL-5.

Estymandem pozostaje asocjacyjny iloraz szans z dwustronnym 95-procentowym przedziałem ufności. Definicje populacji, ekspozycji i wyniku pozostają zgodne z zamrożoną specyfikacją pracy. Deklaracja użycia oraz nasilenie objawów dotyczą pomiaru przekrojowego.

Przed badaniem z udziałem ludzi należy zatwierdzić wymagania etyczne, zasady zgody i ochronę danych.

\section{Plan analizy}

Model regresji logistycznej oszacuje współczynnik dla standaryzowanego PCL-5 z zestawem zmiennych towarzyszących określonym w planie analizy pracy magisterskiej.

Zamrożenie estymandu, reguł oceny jakości i obsługi braków nastąpi przed dostępem do danych wynikowych. Raport obejmie oszacowanie kontrastu głównego z przedziałem niepewności i pełne mianowniki. Jeśli przewidziano hipotezy drugorzędowe, przed analizą należy ustalić sposób kontroli wielokrotności.

\section{Wyniki}

W niniejszym pliku nie umieszczono wyników empirycznych, statystyk ani nieprzeprowadzonych analiz. Sekcja wymaga uzupełnienia po uzyskaniu i weryfikacji materiału dowodowego.

\section{Dyskusja i ograniczenia}

Ograniczenia obejmują możliwe błędy pomiaru, dobór próby, alternatywne mechanizmy i ograniczoną możliwość uogólniania. Korelacja nie dowodzi przyczynowości; poprawność techniczna nie jest dowodem znaczenia klinicznego.

\section{Wnioski}

Hipotezy i sposób pomiaru stanowią podstawę oceny wykonalności oraz przygotowania właściwego protokołu.

\section*{Bibliografia}

{\setlength{\parindent}{0pt}
\setlength{\parskip}{5pt}
\hangindent=1.27cm
Blevins, C. A., Weathers, F. W., Davis, M. T., Witte, T. K., \& Domino, J. L. (2015). The Posttraumatic Stress Disorder Checklist for DSM-5 (PCL-5): Development and initial psychometric evaluation. \textit{Journal of Traumatic Stress, 28}(6), 489--498. \url{https://doi.org/10.1002/jts.22059}\par
\hangindent=1.27cm
Fitzpatrick, K. K., Darcy, A., \& Vierhile, M. (2017). Delivering cognitive behavior therapy to young adults with symptoms of depression and anxiety using a fully automated conversational agent (Woebot): A randomized controlled trial. \textit{JMIR Mental Health, 4}(2), Article e19. \url{https://doi.org/10.2196/mental.7785}\par
}

\end{document}
